This chapter reviews the literature relevant to the proposed hybrid
FAIRness assessment framework. It first introduces the FAIR principles
and their adaptation to research software, before examining existing
FAIRness assessment approaches, research on software documentation
quality, and LLM-based document analysis. Based on this literature,
the chapter derives the requirements for the proposed framework and
identifies the research gap addressed by this thesis.
\section{The FAIR Principles}

The FAIR principles were introduced by Wilkinson et al. 
\cite{wilkinson2016fair} as a minimal, domain-independent framework 
for improving the findability, accessibility, interoperability, and 
reusability of digital research objects. The principles were designed 
to support both human and machine interaction with research outputs, 
with machine-actionability as a central goal — the ability of 
automated agents to discover, evaluate, and act on digital objects 
without prior knowledge of them.

The four FAIR principles are defined as follows.

A digital object is \textbf{Findable} if it is assigned a globally 
unique and persistent identifier, described with rich metadata, and 
indexed in a searchable resource.

It is \textbf{Accessible} if it can be retrieved via an open and 
standardized protocol, and if its metadata remains accessible even 
when the object itself is no longer available.

\textbf{Interoperability} requires that the object uses formal 
knowledge representation languages, FAIR vocabularies, and qualified 
references to other objects.

Finally, an object is \textbf{Reusable} if it has rich attributes, 
a clear license, detailed provenance, and meets community standards.

An important characteristic of the FAIR principles is that they are 
modular and independent — they can be adopted incrementally and in 
any combination. Furthermore, FAIRness is treated as a continuum 
rather than a binary property: a digital object can be more or less 
FAIR depending on how completely it satisfies each principle. This 
continuum directly motivates the development of automated assessment 
tools capable of measuring FAIRness at scale.

The FAIR principles have become one of the most influential frameworks 
in research data management and open science. They have been adopted 
by funding agencies, research infrastructures, publishers, and 
scientific communities worldwide as a reference framework for 
improving the management and reuse of digital research outputs. Their 
widespread adoption has also motivated the development of FAIR 
assessment tools, metrics, and maturity indicators.

Despite their influence, the FAIR principles intentionally do not 
prescribe specific technologies, implementation strategies, or minimum 
compliance requirements. As a result, determining the degree of 
FAIRness of a digital object can be challenging and often requires 
interpretation. This has led to the development of automated 
assessment frameworks and community-specific implementations that 
translate the principles into measurable indicators.

While the FAIR principles were originally developed for research data, 
Wilkinson et al. explicitly noted their applicability to other 
scholarly digital objects including software, workflows, and 
algorithms. Consequently, while the FAIR principles provide a general 
conceptual framework, additional guidance is required to address 
software-specific characteristics and to enable meaningful FAIRness 
assessment of research software repositories.

\section{FAIR for Research Software}

Although the FAIR principles were originally formulated for research
data, their application to research software requires software-specific
interpretation. Research software is a living and complex research
object whose executability, versioning, dependencies, and continuing
evolution distinguish it from conventional research data
\cite{katz2021fair}. These characteristics make a direct application
of data-oriented FAIR practices insufficient and motivate dedicated
principles for research software.

To address these challenges, the FAIR for Research Software (FAIR4RS) 
principles were introduced by Barker et al. \cite{barker2022fair4rs}, 
building on earlier community discussion of how the FAIR principles 
apply to software \cite{lamprecht2020fairsoftware}.
The FAIR4RS principles were developed through a large community 
consultation process involving approximately 500 contributors from 
more than 110 organisations, reflecting broad international consensus 
on the need for software-specific FAIR guidance. The principles adapt 
the original framework to account for the unique characteristics of 
software, introducing new sub-principles where necessary.

Under FAIR4RS, software is \textbf{Findable} if it is assigned 
globally unique and persistent identifiers, with distinct identifiers 
for different versions and components — an extension not present in 
the original FAIR principles. Software is \textbf{Accessible} if it 
can be retrieved via standardised protocols and its metadata remains 
available even if the software itself becomes unavailable. 
\textbf{Interoperability} in FAIR4RS extends beyond metadata to 
include API interactions and references to related software objects. 
Finally, software is \textbf{Reusable} if it is executable, 
understandable, modifiable, and integrable into other systems, with 
clear licensing, provenance, and community standards.
Software identification and citation are also closely connected to
FAIRness. The Software Citation Principles emphasise software as a
legitimate research product and call for citations that support
credit, attribution, unique identification, persistence, and
accessibility \cite{smith2016softwarecitation}. These principles
reinforce the importance of persistent identifiers and citation
metadata in the FAIR assessment of research software.

Research software plays a critical role in enabling reproducible 
science, and FAIR4RS therefore positions software discoverability, 
accessibility, and reusability as essential requirements for 
transparent and reproducible research workflows.

A key insight of FAIR4RS is that FAIRness for software should begin 
during development rather than only at publication time. Unlike data, 
software may be used directly from its development environment, making 
early adoption of FAIR practices essential for reproducibility and 
reuse. The principles have received growing institutional support from 
organisations such as ELIXIR, ARDC, the Netherlands eScience Center, 
and ZB MED.

Despite this progress, FAIR4RS defines high-level principles without 
prescribing concrete automated assessment methods. Many FAIR4RS 
indicators rely on metadata and structural information, while 
qualitative aspects such as documentation clarity, installation 
guidance, and usage examples remain difficult to evaluate 
automatically. This gap between the principles and their practical 
assessment motivates the development of automated tools capable of 
evaluating both structured and unstructured aspects of research 
software repositories.

\section{Existing FAIRness Assessment Approaches}

Several tools and frameworks have been developed to automate the 
assessment of FAIRness for digital research objects. This section 
reviews the most relevant approaches, highlighting their capabilities 
and limitations in the context of research software assessment.

\subsection{FAIR Assessment Frameworks and Maturity Indicators}

The translation of the FAIR principles into measurable criteria has
been addressed by several general FAIR assessment frameworks.
Wilkinson et al. \cite{wilkinson2019evaluating} proposed a scalable
and automated evaluation framework based on community-defined FAIR
Maturity Indicators and associated compliance tests. The approach
demonstrates how high-level FAIR principles can be operationalised
through explicit, machine-actionable tests while retaining
metric-level assessment results.

The need for consistent operational criteria was further addressed by
the FAIR Data Maturity Model developed by Bahim et al.
\cite{bahim2020fairdmm}. The model provides a common set of indicators
and priorities intended to harmonise FAIR assessment approaches and
improve comparability between evaluations. Together, these works
illustrate an important transition from conceptual FAIR principles
towards explicit and reproducible assessment criteria. However, their
primary focus is FAIR assessment of data and digital research objects
rather than qualitative documentation assessment of research software
repositories.

\subsection{F-UJI}

F-UJI is an automated web-based FAIR assessment tool developed within 
the FAIRsFAIR initiative \cite{fuji2021}. It evaluates the FAIRness 
of digital research objects by performing machine-actionable checks 
against predefined metrics, producing quantitative scores across the 
four FAIR dimensions. F-UJI accepts persistent identifiers and 
metadata endpoints as input and evaluates indicators such as metadata 
presence, persistent identifiers, licensing, semantic vocabularies, 
provenance information, and accessibility protocols.

F-UJI treats FAIRness as a spectrum rather than a binary property, 
providing detailed metric-level reporting and evidence traces. It 
demonstrates that many aspects of FAIR compliance can be assessed 
automatically through metadata checks. However, F-UJI is primarily 
metadata-centric and was originally designed for FAIR digital objects 
and datasets rather than research software repositories specifically. 
It was not designed for direct GitHub repository assessment and 
expects persistent identifiers and structured metadata records as 
input, which many repositories in practice do not provide.

Critically, F-UJI evaluates the existence and machine-readability of 
metadata rather than its semantic quality or usefulness. Documentation 
quality, understandability, and usability are not assessed, creating 
a gap between structural FAIR compliance and meaningful software 
reusability.

\subsection{FAIR-IMPACT Metrics}

The FAIR-IMPACT initiative extended existing FAIR assessment 
approaches by defining 17 concrete, measurable FAIR Research Software 
Metrics (FRSM-01 to FRSM-17) for automated assessment of research 
software against the FAIR4RS principles \cite{chuehong2025frsm}. 
These metrics operationalize FAIR4RS into practical automated checks, 
covering findability, accessibility, interoperability, and 
reusability indicators at the code, project, and repository level.

The metrics use three compliance levels — Essential, Important, and 
Useful — and are designed to be domain-agnostic while adaptable to 
disciplinary contexts. Planned integration into F-UJI demonstrates 
their practical application potential. However, the FAIR-IMPACT 
metrics explicitly acknowledge that automated tools can verify the 
existence of metadata but cannot reliably determine its quality, 
relevance, or correctness. Qualitative aspects such as documentation 
clarity and usability remain outside the scope of automated rule-based 
assessment.

\subsection{FAIR-Aware}

In contrast to automated assessment tools, FAIR-Aware is a 
human-centred self-assessment tool \cite{akerman2021fairaware} developed within the FAIRsFAIR 
initiative that focuses on increasing FAIR literacy among researchers. 
It uses a questionnaire-based approach to help researchers reflect on 
how FAIR practices can be integrated into their workflows, providing 
guidance and recommendations rather than automated compliance metrics.

FAIR-Aware demonstrates that FAIRness assessment is not limited to 
verifying machine-readable metadata. Many FAIR considerations require 
researchers to interpret community standards, evaluate reuse 
conditions, and make contextual decisions that depend on domain 
knowledge and qualitative judgment. This highlights the importance of 
qualitative reasoning in FAIR evaluation — a dimension that automated 
rule-based tools cannot fully address.

\subsection{FAIRsharing}

FAIRsharing \cite{sansone2019fairsharing} is a curated registry of FAIR-related standards, 
repositories, databases, policies, and metadata schemas. It provides 
centralised infrastructure to support the discovery, adoption, and 
interconnection of community standards and FAIR-enabling resources 
across research domains. FAIRsharing demonstrates that FAIRness is 
deeply connected to the existence and adoption of shared standards 
and controlled vocabularies — automated FAIR assessment systems rely 
heavily on this ecosystem infrastructure to validate interoperability 
and machine-readability. FAIRsharing is particularly relevant for 
research software because it supports the discovery of community 
standards, metadata schemas, and domain-specific resources that 
contribute to software interoperability and reuse.

\subsection{howfairis}

howfairis \cite{spaaks2020howfairis} is a command-line tool developed
by the Netherlands eScience Center that evaluates a GitHub or GitLab
repository's compliance with the five fair-software.eu
recommendations: whether the repository is publicly accessible with
version control, whether a license file is present, whether the
software is registered in a community registry, whether citation information is available, for example through a
\texttt{CITATION.cff} file — a specification providing human- and
machine-readable citation metadata that developers can include
directly in a repository \cite{druskat2021cff} — and whether quality
checklist evidence is present.
Unlike F-UJI, howfairis operates directly on software repositories rather
than requiring a persistent identifier or metadata endpoint as input,
making it directly applicable to GitHub-hosted research software.

Like F-UJI and the FAIR-IMPACT metrics, howfairis restricts its
assessment to the presence of structural, machine-detectable evidence.
It does not evaluate the content or clarity of README documentation,
the completeness of installation or usage instructions, or any other
qualitative aspect of software reusability.

\subsection{Summary}

Existing FAIRness assessment tools provide valuable automated 
evaluation capabilities and have significantly advanced the practical 
application of FAIR principles. However, their primary focus remains 
on structured metadata and objectively measurable indicators, leaving 
qualitative aspects of software documentation largely unexplored.

\section{Software Documentation Quality}

While existing FAIRness assessment tools focus on structured metadata 
and file presence, research on software documentation quality 
highlights a complementary dimension of software reusability that 
these tools do not address. README files commonly act as an entry
point to software documentation and provide information needed to
understand and use a project. Guidance for scientific software
documentation recommends providing information such as project
purpose, installation and configuration instructions, usage guidance,
licensing information, and links to further documentation
\cite{lee2018documentation}. More broadly, robust research software
also depends on practices that make software easier for users outside
the original development environment to run and reproduce
\cite{taschuk2017robust}.

\subsection{README Content and Structure}

Prana et al. \cite{prana2019readme} conducted the first large-scale 
empirical study of GitHub README files, categorising their content 
into eight types: What, Why, How, When, Who, References, Contribution, 
and Other. Their study of 393 repositories found that while 
installation and usage instructions (How) are present in the majority 
of files, purpose and motivation sections (Why) appear in only 25.7\% 
of repositories, and project status information (When) in only 21.4\%.

These gaps in purpose and status information are not merely
documentation style issues; they can affect a user's ability to
understand the software and determine whether it is suitable for
reuse, making them relevant to FAIR4RS reusability. Rule-based tools can detect whether 
a README file exists, but cannot determine whether it contains the 
information a user needs to evaluate or reuse the software.

\subsection{Beyond Accuracy: Multi-Dimensional Documentation Quality}

Treude et al. \cite{treude2020beyond} proposed a ten-dimensional
framework for assessing software documentation quality, adapted from
the information quality literature. Rather than focusing only on
accuracy and completeness, the framework considers dimensions such as
clarity, understandability, structure, readability, and conciseness.
These dimensions illustrate that the presence or completeness of
documentation alone does not determine its usefulness: a README may
contain the expected information while still being difficult to
understand or use. This distinction motivates the inclusion of
qualitative documentation characteristics alongside structural
FAIRness checks.

\subsection{Practitioner Perspective on Documentation Quality}

Imani et al. \cite{imani2024documentation} provided empirical 
evidence from an industry perspective, finding that over 91\% of 
practitioners rely on OSS documentation for software adoption 
decisions. Their study identified that the top reported challenge is 
not the absence of documentation but the difficulty of finding 
relevant information within it. Nine adoption criteria emerged from 
practitioner interviews, including ease of use, compatibility, 
functionality, and presence of usage examples — criteria that closely 
map to FAIR4RS reusability indicators.

This practitioner evidence reinforces the need for assessment methods 
that evaluate not just the presence but the quality and completeness 
of documentation content.

These findings suggest that structured rubrics and explicit scoring 
criteria are important design considerations for LLM-based assessment 
systems, particularly when consistency and reliability are required.
This picture is corroborated by large-scale empirical work on
documentation quality more broadly. Aghajani et al. \cite{aghajani2019issues}
mined 878 documentation-related artifacts from GitHub issues, pull
requests, mailing lists, and Stack Overflow to construct a taxonomy of
recurring documentation problems, while a companion survey of 146
practitioners \cite{aghajani2020perspective} found persistent gaps
between the documentation developers consider important and what is
actually produced.

The importance of documentation is particularly evident in research
software specifically. Hermann and Fehr \cite{hermann2022documenting}
identify good documentation as a prerequisite for software reuse and
show that research-software documentation practices remain incomplete
and fragmented in practice. Together, these findings support treating
documentation content as an explicit component of software
reusability rather than merely checking whether a README file exists.
\subsection{Summary}

Research on software documentation quality consistently demonstrates 
that documentation presence alone is insufficient to ensure software 
reusability. The content, clarity, and completeness of documentation 
are critical factors that current automated FAIRness assessment tools 
do not evaluate. This creates a clear motivation for approaches that combine 
traditional rule-based indicators with qualitative assessment of 
documentation content.

\section{LLM-Based Document Analysis}

Recent advances in large language models (LLMs) have opened new 
possibilities for automated analysis of natural language artifacts 
in software repositories. This section reviews work on LLM 
capabilities for software engineering tasks, code documentation 
generation, and automated evaluation — providing the technical 
foundation for the AI-based component of this thesis.

\subsection{LLM-Based FAIR Assessment}

A closely related contribution is Fair-Way \cite{sharma2025fairway},
an open-source framework that applies Large Language Models (LLMs) to
automate FAIRness assessment of research datasets. Unlike traditional
rule-based tools, Fair-Way uses LLMs to evaluate FAIRsFAIR Data Object
Assessment Metrics (DOAM), aiming to improve flexibility across
heterogeneous metadata formats and domains.

Fair-Way follows a divide-and-conquer strategy, decomposing
the assessment into fine-grained tasks while splitting large metadata
files into manageable chunks. The system employs prompt engineering
techniques including structured JSON output, task decomposition,
chain-of-thought prompting, and few-shot examples to improve assessment
reliability and reproducibility. The authors evaluate multiple LLMs,
reporting GPT-4o as the best-performing model, and compare Fair-Way
against F-UJI on three datasets from different scientific domains.
Their results demonstrate performance comparable to, and in several
assessment metrics exceeding, the rule-based baseline.

Despite these promising results, Fair-Way differs substantially from
the approach proposed in this thesis. First, it focuses on research
datasets and FAIRsFAIR DOAM metrics, whereas this work targets research
software repositories and FAIR4RS-inspired documentation indicators.
Second, Fair-Way primarily uses LLM-based assessment for the evaluated FAIR
metrics, whereas the approach in this thesis explicitly separates
deterministic repository checks from LLM-based qualitative
documentation assessment. Finally, Fair-Way evaluates dataset metadata, whereas
this thesis combines rule-based repository analysis with LLM-based
assessment of GitHub repository documentation to produce a hybrid
FAIRness evaluation for research software.

These differences highlight an important gap in the current literature.
While Fair-Way demonstrates the potential of LLMs for automated FAIR
assessment, the literature reviewed in this thesis does not identify
an approach that combines rule-based repository analysis with
LLM-based documentation assessment specifically for research software
repositories.

\subsection{LLMs for Software Engineering}

Hou et al. \cite{hou2024llmse} conducted a systematic literature 
review of LLM applications across the software engineering lifecycle, 
finding that LLMs demonstrate strong performance on tasks involving 
natural language artifacts such as documentation, code summarization, 
and explanation generation. A central argument of the review is that 
many software engineering tasks involving documentation, explanation, 
and comprehension are fundamentally language understanding problems 
that benefit from the semantic reasoning capabilities of LLMs.

The review also identifies a consistent recommendation across 
studies: hybrid approaches combining LLMs with traditional methods 
are preferable for reliability-critical applications, as LLM 
limitations such as hallucination, prompt sensitivity, and 
inconsistent outputs warrant complementary rule-based checks. This 
directly supports the hybrid architecture proposed in this thesis.

\subsection{LLM-Based Code Documentation Assessment}

Dvivedi et al. \cite{dvivedi2024llm} conducted a comparative study 
of multiple LLMs for code documentation generation, evaluating 
outputs on six quality dimensions including accuracy, completeness, 
relevance, understandability, and readability. Their study found that 
closed-source models consistently outperform open-source alternatives, 
and that file-level documentation — analogous to README files — is 
consistently the most challenging documentation type for LLMs to 
assess reliably.

The study also validated that checklist-based evaluation rubrics 
reduce inter-rater subjectivity, achieving only 11.67\% disagreement 
across raters. This finding supports the use of structured rubrics 
in LLM-based scoring — the approach adopted in this thesis for 
evaluating README documentation quality against FAIR4RS indicators.

\subsection{LLM-as-Judge}

Zheng et al. \cite{zheng2023llmjudge} conducted the first systematic 
empirical study of the LLM-as-a-judge paradigm, evaluating the use 
of strong LLMs as automated evaluators of open-ended responses. They found that GPT-4 as judge achieves over 80\% agreement with human
raters, approaching human--human agreement levels in the evaluated
settings. Their results demonstrate that strong LLMs can achieve
substantial agreement with human raters in some open-ended evaluation
settings, while also showing that reliability depends on the
evaluation setup and model used.

The study also characterised key biases of LLM judges, including 
position bias, verbosity bias, and self-enhancement bias. For single 
answer grading — the scenario most relevant to README quality 
assessment — the study found that absolute scores are less stable 
than relative comparisons, and that providing structured rubrics with 
explicit criteria substantially improves scoring reliability. These findings motivate the use of explicit scoring criteria and the
investigation of potential scoring biases in the LLM-based assessment
developed in this thesis.
Other work has explored structured LLM-based evaluation strategies.
G-Eval uses explicit evaluation criteria together with a structured
form-filling procedure and reports improved correlation with human
judgements for natural-language generation evaluation
\cite{liu2023geval}. Prometheus similarly investigates fine-grained
LLM evaluation using task-specific score rubrics and reference
materials \cite{kim2024prometheus}. Liu et al. propose AutoCalibrate,
a method that automatically drafts and refines scoring criteria for
an LLM-based evaluator using a small set of human-labelled examples,
showing that calibrating the scoring criteria themselves --- rather
than only the model's output --- can improve alignment with expert
evaluation \cite{liu2024calibrating}. These approaches provide further
motivation for using explicit, dimension-specific scoring rubrics
rather than unconstrained LLM judgements.

Beyond judge-style evaluation, prior work has also investigated
whether language models can reliably assess their own confidence.
Kadavath et al. \cite{kadavath2022know} show that language models can
provide informative estimates of their own uncertainty in several
question-answering settings, while also demonstrating that calibration
varies across tasks and evaluation conditions. More recent work
reports that post-alignment LLMs often exhibit overconfidence, where
expressed confidence does not accurately reflect the model's
correctness rate, and proposes methods to better calibrate
self-reported confidence against actual accuracy
\cite{zhang2024calibrating}. This motivates treating model-reported
confidence as a quantity that requires empirical validation rather
than assuming that it reliably reflects assessment correctness.

\subsection{Summary}

LLMs offer capabilities for analysing natural language documentation
and assessing qualitative aspects of software artefacts that are
difficult to capture through deterministic presence-based checks.
However, existing work also highlights limitations including output
instability, prompt sensitivity, and domain-specific challenges.
These limitations motivate explicit evaluation of LLM reliability,
scoring consistency, and failure modes when LLMs are used for
qualitative assessment.

\section{Requirement Analysis}

The literature reviewed in this chapter suggests a set of concrete
requirements that an automated FAIRness assessment approach for
research software should satisfy. This section makes these
requirements explicit and evaluates the extent to which existing
tools and the framework proposed in this thesis satisfy each of them.

\subsection{Derived Requirements}

Based on the FAIR4RS principles, existing FAIRness assessment
approaches, and research on software documentation quality, the
following requirements are derived for an automated FAIRness
assessment approach for research software:

\begin{description}

\item[R1 --- Structural FAIR Assessment]
The approach should automatically assess machine-detectable
repository characteristics relevant to FAIRness, such as identifiers,
metadata, licensing, citation information, and other structured
software artefacts.

\item[R2 --- Research Software Applicability]
The assessment should address characteristics specific to research
software rather than being limited to research datasets or generic
digital objects.

\item[R3 --- Repository-Level Applicability]
The approach should be applicable directly to software repositories,
such as GitHub repositories, without requiring a separate dataset
record or metadata landing page as the primary assessment target.

\item[R4 --- Qualitative Documentation Assessment]
The approach should assess the content of unstructured software
documentation, including aspects such as installation, usage,
purpose, dependencies, examples, and contribution information,
rather than merely detecting the presence of documentation.

\item[R5 --- Transparent Assessment Evidence]
Assessment results should provide interpretable evidence or
explanations that allow users to understand how individual results
were obtained.

\item[R6 --- Automated and Scalable Execution]
The assessment should be executable automatically across multiple
repositories without requiring manual evaluation for every software
project.

\item[R7 --- Explicit and Reproducible Criteria]
The assessment should apply predefined and documented criteria so
that the evaluation procedure can be reproduced and the meaning of
individual scores or checks remains clear.

\item[R8 --- Reliability Assessment]
For qualitative assessments involving uncertain judgements, the
reliability of the assessment should be evaluated empirically.
Where confidence information is provided by the assessment method,
its relationship with actual assessment correctness should also be
validated rather than assumed.

\end{description}

\subsection{Comparison Against Existing Systems}

Table~\ref{tab:requirements} summarises the extent to which F-UJI,
the FAIR-IMPACT metrics, FAIR-Aware, howfairis, and Fair-Way satisfy
these requirements, compared with the framework proposed in this
thesis.

\begin{table}[h]
\centering
\caption{Comparison of existing FAIRness assessment approaches against the derived requirements. \checkmark: satisfied; (\checkmark): partially satisfied; --: not satisfied.}
\label{tab:requirements}
\begin{tabular}{lcccccc}
\hline
\textbf{Req.} & \textbf{F-UJI} & \textbf{FAIR-} & \textbf{FAIR-} & \textbf{how-} & \textbf{Fair-} & \textbf{Proposed} \\
 & & \textbf{IMPACT} & \textbf{Aware} & \textbf{fairis} & \textbf{Way} & \textbf{Framework} \\
\hline
R1 & \checkmark & \checkmark & -- & \checkmark & (\checkmark) & \checkmark \\
R2 & -- & \checkmark & -- & \checkmark & -- & \checkmark \\
R3 & -- & (\checkmark) & -- & \checkmark & -- & \checkmark \\
R4 & -- & -- & (\checkmark) & -- & (\checkmark) & \checkmark \\
R5 & \checkmark & \checkmark & \checkmark & \checkmark & \checkmark & \checkmark \\
R6 & \checkmark & \checkmark & -- & \checkmark & \checkmark & \checkmark \\
R7 & \checkmark & \checkmark & (\checkmark) & \checkmark & \checkmark & \checkmark \\
R8 & -- & -- & -- & -- & -- & \checkmark \\
\hline
\end{tabular}
\end{table}

The comparison shows that existing approaches address different
parts of the problem. F-UJI and howfairis provide automated,
rule-based assessment of machine-detectable evidence, while the
FAIR-IMPACT metrics extend FAIR assessment towards research
software-specific characteristics. FAIR-Aware incorporates
qualitative human reflection but does not provide automated
repository assessment. Fair-Way demonstrates that LLMs can support
semantic FAIR assessment of structured and unstructured metadata,
but its scope is research data rather than software repository
documentation.

Consequently, the main unresolved combination is direct,
repository-level research software assessment that preserves
deterministic structural checks while additionally evaluating the
qualitative content of software documentation. The framework proposed in this thesis is designed around this
combination and applies predefined and documented assessment criteria
for both components (R7). The reliability of its LLM-based component
is evaluated empirically against manual annotations using agreement
and error measures. The relationship between model-reported
confidence and assessment correctness is also examined rather than
assumed, satisfying R8.

\section{Research Gap}

The preceding analysis shows that existing FAIRness assessment
approaches address complementary aspects of the assessment problem,
but do not provide the combination investigated in this thesis.
Rule-based approaches offer reproducible assessment of
machine-detectable evidence, while recent LLM-based approaches
demonstrate the potential for semantic assessment of less structured
information \cite{hou2024llmse,sharma2025fairway}.

Within the literature reviewed in this thesis, no approach was
identified that combines deterministic repository-level FAIR checks
with LLM-based assessment of qualitative README documentation
specifically for research software. This leaves two questions open:
whether LLM-based documentation assessment provides useful
information beyond rule-based repository checks, and how reliably an
LLM can perform such qualitative assessment.

This thesis addresses this gap through a hybrid assessment approach
in which the two components serve complementary roles. Deterministic
rules assess structural repository evidence, while an LLM evaluates
predefined qualitative documentation dimensions. Rather than merging
both components into a single FAIRness score, their results are
reported together in a per-repository hybrid profile. The evaluation
therefore examines both the complementary contribution of the
LLM-based component and the reliability and failure modes of its
qualitative judgements.